\section{Q17 forecast-to-decision study: protocol and full results}
\label{app:q17}

\paragraph{Data and features.} Q17 uses $30$ admitted native on-change best-bid/offer
contracts across BTC and ETH over $15$ paired calendar dates. Features are past-only quote
statistics computed on the receipt clock. The target is the sign of a short-horizon
mid-quote move discretized into up/down/flat. Feature construction and normalization are
fit on the training portion only and frozen before any test date is scored. Six model
families produce frozen forecasts: a class-prior baseline, logistic regression, gradient
boosting, a multilayer perceptron, a random control, and an oracle with access to the
realized label.

\paragraph{Policies and utility.} Each frozen forecast passes through three fixed policies.
The \emph{argmax} policy takes the highest-probability direction. The \emph{confidence}
policy acts only above a fixed probability threshold and otherwise abstains. The
\emph{scheduling} policy allocates activity across a $12$-opportunity window. Utility is
computed with inverse-contract price arithmetic and scenario fees at a $5$ bps-per-side
cost and a $100$ ms quote-offset latency. Utility uses quote endpoints and assumed costs. It does not observe realized fills, funding, or endogenous impact, so it is a decision
utility, not a realized profit.

\paragraph{Joint gate.} A forecast-to-decision improvement is declared only when a model
shows both a predictive advantage and a utility reversal across at least two policy families
on six of eight paired dates. Table~\ref{tab:q17primary} reports the four primary contrasts
per asset (one macro-F1 and three utility contrasts). No contrast passes: the BTC macro-F1
gain is significant, but the confidence and scheduling contrasts fail, and all four ETH
contrasts fail. The confidence policy collapses to zero utility (abstain) for every learned
model, and the scheduling contrast is negative on four of eight dates for both assets.

\begin{table}[h]\centering
\caption{Q17 primary contrasts (MLP minus logistic) with bootstrap intervals over eight
paired dates. Source: \texttt{experiments/results/q17\_primary\_contrasts.csv}. No
contrast passes the joint gate.}
\label{tab:q17primary}
\begin{tabular}{llrrrr}
\toprule
asset & metric & mean & 95\% interval & neg. dates & passed \\
\midrule
BTC & macro\_f1 & 0.142103 & [0.039529, 0.244677] & 0 & \checkmark \\
BTC & argmax & 2.122375 & [0.292744, 3.952005] & 0 & $\times$ \\
BTC & confidence & 0.000000 & [0.000000, 0.000000] & 0 & $\times$ \\
BTC & scheduling & 0.025961 & [-0.102278, 0.154200] & 4 & $\times$ \\
ETH & macro\_f1 & 0.097081 & [0.006888, 0.187274] & 1 & $\times$ \\
ETH & argmax & 1.801718 & [0.096671, 3.506765] & 1 & $\times$ \\
ETH & confidence & 0.000000 & [0.000000, 0.000000] & 0 & $\times$ \\
ETH & scheduling & -0.070943 & [-0.524662, 0.382776] & 4 & $\times$ \\
\bottomrule
\end{tabular}
\end{table}


\begin{table}[h]\centering
\caption{Q17 absolute performance by model family. Utility in basis points per
opportunity (argmax, confidence) or per 12-opportunity window (scheduling) at 5 bps per
side. Source: \texttt{experiments/results/q17\_figure\_data.json}.}
\label{tab:q17abs}
\begin{tabular}{llrrrr}
\toprule
asset & model & macro-F1 & argmax & confidence & scheduling \\
\midrule
BTC & prior & 0.167633 & -10.143430 & 0.000000 & -0.087832 \\
BTC & logistic & 0.361085 & -9.741636 & 0.000000 & 0.680911 \\
BTC & hist\_gb & 0.486577 & -8.034500 & 0.000000 & 0.531820 \\
BTC & mlp & 0.503188 & -7.619262 & 0.000000 & 0.706872 \\
BTC & random & 0.318313 & -6.772006 & 0.000000 & -0.080510 \\
BTC & oracle & 1.000000 & -5.673321 & 0.000000 & 1.997775 \\
ETH & prior & 0.171960 & -10.455921 & 0.000000 & -0.108507 \\
ETH & logistic & 0.352168 & -9.956131 & 0.000000 & 0.903621 \\
ETH & hist\_gb & 0.371520 & -9.559819 & 0.000000 & 0.797432 \\
ETH & mlp & 0.449249 & -8.154412 & 0.000000 & 0.832678 \\
ETH & random & 0.319194 & -6.992946 & 0.000000 & -0.080479 \\
ETH & oracle & 1.000000 & -5.461902 & 0.000000 & 2.752498 \\
\bottomrule
\end{tabular}
\end{table}


\begin{table}[h]\centering
\caption{Q17 per-date macro-F1 advantage of the MLP over logistic regression, by paired date,
with the mean. Source: \texttt{experiments/results/q17\_figure\_data.json}.}
\label{tab:q17days}
\begin{tabular}{lrrrrrrrrr}
\toprule
asset & d1 & d2 & d3 & d4 & d5 & d6 & d7 & d8 & mean \\
\midrule
BTC & 0.223477 & 0.072209 & 0.265424 & 0.076891 & 0.089394 & 0.088852 & 0.194926 & 0.125649 & \textbf{0.142103} \\
ETH & -0.001203 & 0.073973 & 0.190915 & 0.060807 & 0.086582 & 0.080443 & 0.196469 & 0.088661 & \textbf{0.097081} \\
\bottomrule
\end{tabular}\end{table}


Table~\ref{tab:q17days} gives the per-date macro-F1 advantage: the BTC gain is positive on all
eight dates (hence the significant contrast), while the ETH gain is negative on one date, which
is why the ETH predictive contrast is weaker even though its mean is positive.

\paragraph{Absolute performance.} Table~\ref{tab:q17abs} gives the absolute macro-F1 and
per-policy utility for every model family. The ordering by macro-F1 (oracle $>$ mlp $>$
hist\_gb $>$ logistic $>$ random $>$ prior) does not match the ordering by argmax utility,
and every learned argmax utility is negative. The oracle's own argmax utility is negative
($-5.67$ BTC, $-5.46$ ETH), which bounds how much the cost layer alone removes: even a
perfect direction forecast trades at a loss under this policy and cost model. This is the
quantitative core of the section: the forecasting leaderboard and the decision leaderboard
are different orderings, and the gap is large enough to invert the deployment choice.

\paragraph{Sensitivity across later dates.} A frozen follow-up sweeps fixed cost and latency
scenarios on additional dates without changing the gate. The per-policy utility ranges
remain policy-dependent and continue to straddle zero for the confidence and scheduling
families, so the negative is not an artifact of one cost assumption. We report ranges rather
than a single tuned scenario precisely to avoid selecting the cost point that would flip the
conclusion.
